% ============================================================================
%  presentation/appscreens.tex
%
%  The Android application, drawn.
%
%  WHAT THESE ARE. Vector mock-ups of four screens, drawn in TikZ from the
%  Compose source in mobile/app/src/main/java/org/noisehanoi/mobile/ui/. Every
%  string, every colour and every control on them is taken from that source:
%
%    Home     <- ui/HomeScreen.kt          Measure  <- ui/FormScreen.kt
%    Map      <- ui/MapScreen.kt           Results  <- ui/ResultsScreen.kt
%
%  Colours are the app's own Material 3 scheme from ui/Theme.kt, and the map
%  bands are levelColour() from the same file.
%
%  WHAT THESE ARE NOT. They are not device captures. Nobody on the team has run
%  the app on a handset with a sound source in front of it — that is stated on
%  the "what the app refuses to claim" slide and it is why there are no real
%  screenshots here yet. Replace these with captures the day the app is run on a
%  real phone: put the PNGs in presentation/screens/ and swap the \phone... call
%  for an \includegraphics. Nothing else in main.tex has to change.
%
%  Each macro draws one phone at a nominal 44 mm x 92 mm. Wrap in \scalebox to
%  resize; the text scales with it.
% ============================================================================

% ---------------------------------------------------------------- app palette
\definecolor{appink} {HTML}{0B3D5C}   % Theme.kt  Ink   -> primary
\definecolor{appsky} {HTML}{3E8FB8}   % Theme.kt  Sky   -> secondary
\definecolor{appsurf}{HTML}{F4F6F8}   % Material 3 surfaceVariant, light scheme
\definecolor{appline}{HTML}{C9D2D8}
\definecolor{apptext}{HTML}{5A6A75}   % colorScheme.outline, the grey body text

% levelColour() in ui/Theme.kt, band for band.
\definecolor{lvlA}{HTML}{2E7D32}      % < 55
\definecolor{lvlB}{HTML}{9E9D24}      % 55-65
\definecolor{lvlC}{HTML}{EF6C00}      % 65-70
\definecolor{lvlD}{HTML}{D84315}      % 70-80
\definecolor{lvlE}{HTML}{B71C1C}      % > 80

% ---------------------------------------------------------------- primitives
\tikzset{
  phonebody/.style   = {draw=appink!55, line width=.5pt, rounded corners=2.2mm,
                        fill=white},
  phonescreen/.style = {fill=white, draw=none},
  appbar/.style      = {fill=appink, draw=none},
  card/.style        = {draw=appline, line width=.3pt, rounded corners=1mm,
                        fill=white},
  cardfill/.style    = {card, fill=appsurf},
  btn/.style         = {fill=appink, draw=none, rounded corners=1.6mm},
  obtn/.style        = {draw=appink, line width=.3pt, rounded corners=1.6mm,
                        fill=white},
  chip/.style        = {draw=appline, line width=.3pt, rounded corners=1.2mm,
                        fill=white},
  chipon/.style      = {chip, fill=appink!12, draw=appink!45},
  ttl/.style         = {font=\tiny\bfseries, text=black, anchor=north west,
                        inner sep=0pt},
  bdy/.style         = {font=\fontsize{3.4}{4}\selectfont, text=apptext,
                        anchor=north west, inner sep=0pt},
  btnlab/.style      = {font=\fontsize{3.6}{4}\selectfont, text=white,
                        anchor=center, inner sep=0pt},
  obtnlab/.style     = {font=\fontsize{3.6}{4}\selectfont, text=appink,
                        anchor=center, inner sep=0pt},
}

% \phoneframe{<body>}{<app bar title>} — chrome shared by all four screens.
% Draws the handset, the status bar and the Material 3 top app bar, then places
% <body> in a scope whose origin is the top-left of the content area.
\newcommand{\phoneframe}[2]{%
  \begin{tikzpicture}[x=1mm, y=1mm, line cap=round, font=\sffamily]
    \sffamily
    \draw[phonebody] (-1.6,2.2) rectangle (45.6,-93.6);
    \fill[phonescreen] (0,0) rectangle (44,-92);
    % status bar
    \node[font=\fontsize{2.8}{3}\selectfont, text=apptext, anchor=north west]
      at (1.8,-0.4) {09:41};
    \node[font=\fontsize{2.8}{3}\selectfont, text=apptext, anchor=north east]
      at (42.2,-0.4) {LTE \raisebox{-.1ex}{\rule{1.6mm}{1.1mm}}};
    % top app bar
    \fill[appbar] (0,-3.4) rectangle (44,-11.4);
    \node[font=\scriptsize, text=white, anchor=west] at (2.4,-7.4) {#2};
    % Content is clipped to the screen exactly as it is on the device: these
    % screens scroll, and a mock-up that spills past the bezel reads as a bug.
    \begin{scope}
      \clip (0,-3.4) rectangle (44,-92);
      \begin{scope}[shift={(0,-11.4)}]
        #1
      \end{scope}
    \end{scope}
  \end{tikzpicture}%
}

% \appcard{<y top>}{<height>}{<inner>} — a Material 3 Card. Inner content is
% placed from the card's top-left plus the 2 mm padding the app uses.
\newcommand{\appcard}[4][card]{%
  \draw[#1] (2,#2) rectangle (42,#2-#3);
  \begin{scope}[shift={(4,#2-2)}] #4 \end{scope}%
}

% \appbutton{x}{y}{w}{label}
\newcommand{\appbutton}[4]{%
  \draw[btn] (#1,#2) rectangle (#1+#3,#2-4.6);
  \node[btnlab] at (#1+#3/2,#2-2.3) {#4};}
\newcommand{\appobutton}[4]{%
  \draw[obtn] (#1,#2) rectangle (#1+#3,#2-4.6);
  \node[obtnlab] at (#1+#3/2,#2-2.3) {#4};}

% Grey text ruled as texture: body copy at this size is a shape, not a read.
\newcommand{\bodylines}[3]{%
  \foreach \i/\w in {#3} {\fill[apptext!35] (#1,#2-\i*2.2) rectangle
                          (#1+\w,#2-\i*2.2-1.1);}}

% ============================================================================
%  1. Home  --  ui/HomeScreen.kt
% ============================================================================
\newcommand{\phoneHome}{%
\phoneframe{%
  \node[bdy, text width=38mm] at (3,-2)
    {Field collection for the Hanoi urban noise campaign. Forms are the ones
     deployed on KoboToolbox.};

  \node[ttl] at (3,-11) {Collect};

  \appcard{-13.5}{16}{%
    \node[ttl] at (0,0) {Noise measurement};
    \node[bdy] at (0,-3) {16 questions {\textperiodcentered} v3};
    \appbutton{0}{-6}{18}{Fill this form}}

  \appcard{-31}{16}{%
    \node[ttl] at (0,0) {Construction site};
    \node[bdy] at (0,-3) {6 questions {\textperiodcentered} v1};
    \appbutton{0}{-6}{18}{Fill this form}}

  \node[ttl] at (3,-49) {The study};

  \appcard{-51.5}{13}{%
    \node[ttl] at (0,0) {Predicted noise map};
    \node[bdy, text width=34mm] at (0,-3)
      {The 40\,m grid, by hour, with the 363 points.};
    \appobutton{0}{-7.4}{16}{Open the map}}

  \appcard{-66}{13}{%
    \node[ttl] at (0,0) {What the study found};
    \node[bdy, text width=34mm] at (0,-3)
      {The delivered model and the three negative results.};
    \appobutton{0}{-7.4}{17}{Open the results}}
}{Noise Hanoi}}

% ============================================================================
%  2. Measure  --  ui/FormScreen.kt, the measurement card
% ============================================================================
\newcommand{\phoneMeasure}{%
\phoneframe{%
  \node[bdy] at (3,-2) {Question 12 of 16};
  \bodylines{3}{-5}{0/26, 1/34, 2/18}

  \appcard[cardfill]{-14}{46}{%
    \node[ttl] at (0,0) {Measurement and audio clip};
    \node[bdy, text width=34mm] at (0,-3.2)
      {One microphone session: 25\,s, A-weighted, SLOW. The clip submitted is
       the stretch that was measured.};
    % the live reading, in levelColour()
    \node[font=\fontsize{11}{11}\selectfont, text=lvlD, anchor=north west]
      at (0,-14) {71.4 dB};
    \node[bdy] at (0,-22.5) {L\textsubscript{eq} 70.8 dB over 18 s};
    % LinearProgressIndicator
    \fill[appink!18, rounded corners=.4mm] (0,-26) rectangle (34,-27);
    \fill[appink,    rounded corners=.4mm] (0,-26) rectangle (24.5,-27);
    \node[bdy] at (0,-29) {0.00 \% of samples at full scale};
    \appbutton{0}{-32}{15}{Measure 25 s}
    \appobutton{16.5}{-32}{9}{Stop}}

  \appcard{-62}{11}{%
    \node[ttl] at (0,0) {Position};
    \node[bdy] at (0,-3) {\textpm\,6\,m {\textperiodcentered} inside the 10\,m gate};
    \node[bdy] at (0,-6) {20.99220, 105.94426};}

  \appbutton{2}{-73.5}{40}{Save to outbox and send}
  \node[bdy] at (3,-79.8) {Offline: 2 waiting in the outbox.};
}{Noise measurement}}

% ============================================================================
%  3. Map  --  ui/MapScreen.kt
% ============================================================================
\newcommand{\phoneMap}{%
\phoneframe{%
  % FilterChip row over the three sites
  \draw[chipon] (2,-1.5) rectangle (16,-6);
    \node[font=\fontsize{3.2}{3.6}\selectfont, text=appink] at (9,-3.75) {Ocean Park};
  \draw[chip] (17,-1.5) rectangle (30,-6);
    \node[font=\fontsize{3.2}{3.6}\selectfont, text=apptext] at (23.5,-3.75) {Hoan Kiem};
  \draw[chip] (31,-1.5) rectangle (42,-6);
    \node[font=\fontsize{3.2}{3.6}\selectfont, text=apptext] at (36.5,-3.75) {Vinh Tuy};

  % NoiseCanvas: the 40 m grid. Deterministic pseudo-random bands so the
  % picture is stable between builds -- it is a mock-up, not data.
  \begin{scope}
    \clip (2,-8) rectangle (42,-45);
    % A smooth field quantised to the app's bands. The kernel keys on distance
    % to a road, so the picture has to fall off away from the roads rather than
    % scatter: this is a sum of low-frequency terms, banded the way the app
    % bands its own output.
    \foreach \r in {0,...,17} {
      \foreach \c in {0,...,19} {
        \pgfmathsetmacro{\v}{1.55 + 1.05*sin((\c*17+22)) + 0.75*cos((\r*23-14))
                              + 0.55*sin(((\r+\c)*11)) + 0.35*cos(((\r-\c)*29))}
        \pgfmathtruncatemacro{\band}{max(0, min(3, int(\v)))}
        \edef\bc{\ifcase\band lvlA\or lvlB\or lvlC\or lvlD\fi}
        \fill[\bc, opacity=.9] (2+\c*2,-8-\r*2.06)
                               rectangle ++(2,-2.06);}}
    % The two arterials, with the one-cell falloff either side.
    \fill[lvlC] (2,-19.4) rectangle (42,-25);
    \fill[lvlD] (2,-21) rectangle (42,-23.4);
    \fill[lvlC] (24.4,-8) rectangle (30,-45);
    \fill[lvlD] (26,-8) rectangle (28.4,-45);
    % the 363 measured points, when "Show the 363 points" is on
    \foreach \x/\y in {8/-14, 13/-19, 19/-27, 27/-13, 31/-31, 22/-38, 35/-22,
                       11/-33, 38/-36, 16/-11} {
      \fill[white] (\x,\y) circle (.75mm);
      \fill[appink] (\x,\y) circle (.45mm);}
  \end{scope}
  \draw[appline, line width=.3pt] (2,-8) rectangle (42,-45);

  % Legend()
  \foreach \i/\c/\l in {0/lvlA/{<55}, 1/lvlB/{55-65}, 2/lvlC/{65-70},
                        3/lvlD/{70-80}, 4/lvlE/{>80}} {
    \fill[\c] (2+\i*8,-47) rectangle ++(2,-2);
    \node[font=\fontsize{2.8}{3}\selectfont, text=apptext, anchor=west]
      at (4.4+\i*8,-48) {\l};}

  \node[ttl] at (2,-52) {Hour: 17:00};
  \fill[appink!20, rounded corners=.4mm] (2,-56) rectangle (42,-56.9);
  \fill[appink,    rounded corners=.4mm] (2,-56) rectangle (31,-56.9);
  \fill[appink] (31,-56.45) circle (1.1mm);

  \node[ttl] at (2,-59.5) {Traffic volume: x1.00 of measured};
  \fill[appink!20, rounded corners=.4mm] (2,-63.5) rectangle (42,-64.4);
  \fill[appink,    rounded corners=.4mm] (2,-63.5) rectangle (16,-64.4);
  \fill[appink] (16,-63.95) circle (1.1mm);

  \appobutton{2}{-67}{18}{Reset to measured}
  \appobutton{21.5}{-67}{20.5}{Hide the 363 points}

  \appcard{-73.5}{13}{%
    \node[ttl] at (0,0) {Selected cell};
    \node[bdy] at (0,-3) {20.99180, 105.94402};
    \node[bdy] at (0,-6) {68.4 dB predicted at 17:00};}
}{Predicted noise map}}

% ============================================================================
%  4. Results  --  ui/ResultsScreen.kt
% ============================================================================
\newcommand{\phoneResults}{%
\phoneframe{%
  \appcard[cardfill]{-2}{22}{%
    \node[ttl] at (0,0) {The campaign};
    \node[bdy] at (0,-3.2) {363 measurements, 2026-06-10 to};
    \node[bdy] at (0,-6.0) {2026-07-22, 47 to 88 dB.};
    \node[bdy] at (0,-9.2) {{\textperiodcentered} Ocean Park --- 184 points};
    \node[bdy] at (0,-11.8) {{\textperiodcentered} Hoan Kiem lake --- 99 points};
    \node[bdy] at (0,-14.4) {{\textperiodcentered} Vinh Tuy area --- 80 points};}

  \node[ttl, text width=38mm] at (2,-26)
    {The three negative results are the contribution};

  \appcard{-33}{13}{%
    \node[ttl, text width=34mm] at (0,0)
      {1. A three-parameter physical law beat every learned model};
    \bodylines{0}{-6}{0/34, 1/34, 2/21}}

  \appcard{-47.5}{11}{%
    \node[ttl, text width=34mm] at (0,0)
      {2. Morphology from OpenStreetMap added nothing measurable};
    \bodylines{0}{-6}{0/34, 1/26}}

  \appcard{-60}{11}{%
    \node[ttl, text width=34mm] at (0,0)
      {3. A model trained in one city did not transfer};
    \bodylines{0}{-4}{0/34, 1/29}}

  \node[ttl] at (2,-70) {Delivered model: physical kernel};
  \node[bdy] at (2,-73.2) {Buffered LOO 300 m};
  % the ProtocolTable header and its first row
  \node[bdy] at (2,-76.4) {Model};
  \node[bdy] at (28,-76.4) {R\textsuperscript{2}};
  \node[bdy] at (36,-76.4) {MAE};
  \draw[appline, line width=.3pt] (2,-77.9) -- (42,-77.9);
  \node[bdy] at (2,-78.8) {physical};
  \node[bdy] at (28,-78.8) {0.246};
  \node[bdy] at (36,-78.8) {5.01};
}{What the study found}}
